\documentclass[12pt]{article}
\usepackage[utf8]{inputenc}
\usepackage[T1]{fontenc}
\usepackage[margin=1in]{geometry}
\usepackage{setspace}
\usepackage{xcolor}
\usepackage[colorlinks=true, urlcolor=blue!60!black, linkcolor=blue!60!black]{hyperref}
\setlength{\parindent}{2em}
\setlength{\parskip}{1.25em}
\renewcommand{\baselinestretch}{1.0}
\begin{document}
\begin{center}
{\LARGE \textit{Loving v.\ Virginia}, 388 U.S.\ 1 (1967)\par}
\vspace{0.8em}
{\large Case Brief --- Four-Page Edition\par}
\end{center}
\vspace{1em}

\section{Memory Jogger}

A unanimous Court struck down Virginia's ban on interracial marriage, holding that racial restrictions on marriage violate both the Equal Protection Clause and the liberty protected by the Due Process Clause of the Fourteenth Amendment (\href{https://supreme.justia.com/cases/federal/us/388/1/}{Justia}).

\section{Detailed Case Facts}

Mildred Jeter, a Black woman, and Richard Loving, a white man, residents of Caroline County, Virginia, married in Washington, D.C., on June 2, 1958, and returned home. In July 1958 police raided their home at night and arrested them. A grand jury indicted them under Virginia's anti-miscegenation scheme, dating to the Racial Integrity Act of 1924, which barred a ``white person and colored person'' from leaving the State to evade the ban and returning to cohabit, and made interracial marriage a felony punishable by one to five years in the penitentiary. Related provisions voided such marriages automatically and defined ``white person'' as having no ascertainable non-Caucasian blood, with an exception for those with one-sixteenth or less American Indian blood. Virginia was one of sixteen States that still banned interracial marriage.

\section{Procedural History}

In January 1959 the Lovings pleaded guilty. Judge Leon Bazile sentenced each to one year in jail, suspended for twenty-five years on condition that they leave Virginia and not return together, writing that ``Almighty God created the races white, black, yellow, malay and red, and he placed them on separate continents.'' From the District of Columbia, they moved in 1963 to vacate the judgment and later filed a federal class action. The Supreme Court of Appeals of Virginia upheld the statutes, relying on \textit{Naim v.\ Naim}, while striking the banishment condition. The U.S.\ Supreme Court reversed on June 12, 1967.

\section{Judicial Votes}

\textbf{9--0}, reversed.
\begin{itemize}
\item \textbf{Majority:} Chief Justice Warren, joined by Justices Black, Douglas, Clark, Harlan, Brennan, White, and Fortas.
\item \textbf{Concurring in the judgment:} Justice Stewart.
\item \textbf{Dissents:} None.
\end{itemize}

\section{Holding}

Statutes that prohibit or punish marriage on the basis of race violate (1) the Equal Protection Clause, because racial classifications---even if ``equally applied''---are subject to ``the most rigid scrutiny'' and require a permissible objective independent of invidious racial discrimination; and (2) the Due Process Clause, because the freedom to marry is a fundamental liberty that cannot be denied on the basis of racial classifications. The Court repudiated \textit{Pace v.\ Alabama} and, in footnote 11, said the result would be the same ``even assuming an even-handed state purpose to protect the `integrity' of all races.''

\section{Analysis of Opinions}

\textbf{Majority (Warren): equal protection.} Marriage is subject to state police power (\textit{Maynard v.\ Hill}), but that power is bounded by the Fourteenth Amendment. Virginia argued that punishing white and Black spouses alike meant no invidious discrimination and only rational-basis review. Following \textit{McLaughlin v.\ Florida}, the Court held that equal application does not immunize a racial classification from ``the very heavy burden of justification.'' Debates in the Thirty-ninth Congress were ``at best, inconclusive'' and could not override the Amendment's ``clear and central purpose.'' Because Virginia barred only interracial marriages involving white persons, the laws were measures ``designed to maintain White Supremacy.''

\textbf{Majority (Warren): due process.} In two short paragraphs, citing \textit{Skinner v.\ Oklahoma}, the Court called marriage one of the ``basic civil rights of man,'' fundamental to our very existence and survival; denying it on racial grounds ``is surely to deprive all the State's citizens of liberty without due process of law.''

\textbf{Stewart (concurring).} Adhering to his \textit{McLaughlin} view, he treated race-based criminal liability as per se invalid: ``it is simply not possible for a state law to be valid under our Constitution which makes the criminality of an act depend upon the race of the actor.''

\section{Examples: Future Applications}

\begin{enumerate}
\item \textbf{Same-Side: Ethnicity Ban.} A ban on marriage across national origins, applied ``equally'' to preserve cultural heritage, fails the most rigid scrutiny.
\item \textbf{Same-Side: Prison Marriages.} A rule letting inmates marry only for ``compelling reasons'' arbitrarily burdens the fundamental right (as in \textit{Turner v.\ Safley}, 1987).
\item \textbf{Opposite-Side: Eligibility Rules.} Licensing, minimum-age, and consanguinity rules are reasonable police-power regulation resting on no suspect classification.
\item \textbf{Opposite-Side: Fees and Waiting Periods.} A modest fee with an indigency waiver and a 72-hour wait do not substantially interfere with the decision to marry.
\item \textbf{Fence-Sitter: Child-Support Precondition.} Barring marriage for parents behind on support involves no racial classification and serves child welfare, yet directly burdens the right to marry; \textit{Zablocki v.\ Redhail} (1978) struck it down.
\end{enumerate}

\section{Critique}

\textbf{The originalist test case.} Critics of originalism, including Richard Posner, Cass Sunstein, and Michael Klarman, note that many framers and ratifiers expected the Fourteenth Amendment to leave anti-miscegenation laws intact. Originalists such as Steven Calabresi and Andrea Matthews respond that the Amendment constitutionalized the Civil Rights Act of 1866's equal right to contract, including marriage, and barred caste legislation. Textualist skeptics of substantive due process accept the equal protection holding while questioning Part~II.

\textbf{Liberty and equality.} Progressive scholars praise \textit{Loving}'s synergy of liberty and equality but find its due process analysis under-theorized, leaving later courts to decide whether the right to marry extends to same-sex couples, as \textit{Obergefell} held. Critical race scholars stress its anti-subordination core: the Court named ``White Supremacy'' as an unconstitutional objective.

\section{Key Quotations}

\begin{quote}
``The freedom to marry has long been recognized as one of the vital personal rights essential to the orderly pursuit of happiness by free men.''
\end{quote}
\noindent --- Chief Justice Warren, for the Court, 388 U.S.\ at 12.

\begin{quote}
``Under our Constitution, the freedom to marry, or not marry, a person of another race resides with the individual, and cannot be infringed by the State.''
\end{quote}
\noindent --- Chief Justice Warren, for the Court, 388 U.S.\ at 12.

\begin{quote}
``At the very least, the Equal Protection Clause demands that racial classifications, especially suspect in criminal statutes, be subjected to the `most rigid scrutiny,' \textit{Korematsu v.\ United States}, 323 U.S.\ 214, 216 (1944), and, if they are ever to be upheld, they must be shown to be necessary to the accomplishment of some permissible state objective, independent of the racial discrimination which it was the object of the Fourteenth Amendment to eliminate.''
\end{quote}
\noindent --- Chief Justice Warren, for the Court, 388 U.S.\ at 11.

\end{document}
